\section{Cluster algebras for SG Y-systems}
\label{sec:clusterSG}

In this section we identify $\mathbb{T}_{n}(\mathrm{SG})$ and
$\mathbb{Y}_{n}(\mathrm{SG})$ as relations for cluster variables
and coef\/f\/icients of the cluster algebra
associated with a certain quiver $Q_{n}(\mathrm{SG})$.
We follow \cite{Fomin07}
 for def\/initions and conventions concerning cluster algebras
with coef\/f\/icients,
which are summarized in Appendix \ref{sec:groc} for the reader's
convenience.

\subsection{Parity decompositions of T and Y-systems}


For $(i,u)\in \mathcal{I}_{n}$,
we set the parity condition $\mathbf{P}_{+}$ by,
for even $n$,
\begin{gather*}
%\label{eq:GPcond1}
\mathbf{P}_{+}: \
\begin{cases}
\mbox{$i+u$ is even},& i=1,\dots,n,\\
\mbox{$n+u$ is even}, & i=n+1,
\end{cases}
\end{gather*}
and, for odd $n$,
\begin{gather*}
%\label{eq:GPcond2}
\mathbf{P}_{+}: \
\begin{cases}
\mbox{$u$ is even}, & i=1,\\
\mbox{$i+u$ is even}, & i=2,\dots,n, \\
\mbox{$n+u$ is even}, & i=n+1.\\
\end{cases}
\end{gather*}
Let $\mathbf{P}_-$ be the negation of $\mathbf{P}_+$.
We write, for example, $(i,u):\mathbf{P}_{+}$ if $(i,u)\in
\mathcal{I}_n$ satisf\/ies
$\mathbf{P}_{+}$.
Let $\mathcal{I}_{n\varepsilon}$ ($\varepsilon=\pm$)
be the set of all $(i,u):\mathbf{P}_{\varepsilon}$.
Def\/ine $\EuScript{T}^{\circ}_{n}(\mathrm{SG})_{\varepsilon}$
($\varepsilon=\pm$)
to be the subring of $\EuScript{T}^{\circ}_{n}(\mathrm{SG})$
generated by
 $T_i(u)$
$((i,u)\in \mathcal{I}_{n\varepsilon})$.
Then, we have
$\EuScript{T}^{\circ}_{n}(\mathrm{SG})_+
\simeq
\EuScript{T}^{\circ}_{n}(\mathrm{SG})_-
$
by $T_i(u)\mapsto T_i(u+1)$ and
\begin{gather*}
%\label{eq:Tfact}
\EuScript{T}^{\circ}_{n}(\mathrm{SG})
\simeq
\EuScript{T}^{\circ}_{n}(\mathrm{SG})_+
\otimes_{\mathbb{Z}}
\EuScript{T}^{\circ}_{n}(\mathrm{SG})_-.
\end{gather*}

For  $(i,u)\in \mathcal{I}_{n}$,
we introduce another parity condition $\mathbf{P}'_{+}$ by
\begin{gather*}
%\label{eq:GP'cond1}
\mathbf{P}'_{+}: \
\begin{cases}
\mbox{$i+u$ is odd},& i=1,\dots,n,\\
\mbox{$n+u$ is odd}, & i=n+1,
\end{cases}
\end{gather*}
We have
\[
(i,u):\ \mathbf{P}'_+ \ \Longleftrightarrow\
  (i,u\pm d_i): \ \mathbf{P}_+,
\]
where $d_i$ is given in \eqref{eq:di}.
Let $\mathbf{P}'_-$ be the negation of $\mathbf{P}'_+$.
Let $\mathcal{I}'_{n\varepsilon}$
($\varepsilon=\pm$)  be
 the set of all $(i,u):\mathbf{P}'_{\varepsilon}$.
Def\/ine $\EuScript{Y}^{\circ}_{n}(\mathrm{SG})_{\varepsilon}$
($\varepsilon=\pm$)
to be the subgroup of $\EuScript{Y}^{\circ}_{n}(\mathrm{SG})$
generated by
$Y_i(u)$, $1+Y_i(u)$
$((i,u)\in \mathcal{I}'_{n\varepsilon})$.
Then, we have
$\EuScript{Y}^{\circ}_{n}(\mathrm{SG})_+
\simeq
\EuScript{Y}^{\circ}_{n}(\mathrm{SG})_-
$
by $Y_i(u)\mapsto Y_i(u+1)$,
$1+Y_i(u)\mapsto 1+Y_i(u+1)$,
 and
\begin{gather*}
%\label{eq:Yfact}
\EuScript{Y}^{\circ}_{n}(\mathrm{SG})
\simeq
\EuScript{Y}^{\circ}_{n}(\mathrm{SG})_+
\times
\EuScript{Y}^{\circ}_{n}(\mathrm{SG})_-.
\end{gather*}

From now on, we mainly treat the $+$ parts,
$\EuScript{T}^{\circ}_{n}(\mathrm{SG})_+$
and
$\EuScript{Y}^{\circ}_{n}(\mathrm{SG})_+$.

\subsection[Quiver $Q_n(\mathrm{SG})$]{Quiver $\boldsymbol{Q_n(\mathrm{SG})}$}
\label{subsect:quiver}


%%%%%%%%%%%%%%%%%%%%%%%%%%%%%
% SG
\begin{figure}
\centerline{\includegraphics{Nakanishi-Fig2}}
\caption{The quiver $Q_{n}(\mathrm{SG})$  for $n=8$
(upper) and for $n=7$ (lower),
where, except for the leftmost vertex of each quiver $Q_i$,
 all the  vertices in the same position
 in  $n-1$ quivers $Q_1,\dots,Q_{n-1}$ are
identif\/ied.}
\label{fig:quiverG}
\end{figure}

Recall that a {\em quiver\/} is an oriented graph,
namely, it consists of the vertices and the arrows
connecting them.
With each $n \geq 4$ we associate
 a quiver $Q_{n}(\mathrm{SG})$ as below.
First, as rather general examples, the cases $n=8$ and $n=7$ are given in
 Fig.~\ref{fig:quiverG},
where, except for the leftmost vertex of each quiver $Q_i$,
 all the  vertices in the same position
 in  $n-1$ quivers $Q_1,\dots,Q_{n-1}$ are
identif\/ied.
For a general $n$,  the quiver $Q_{n}(\mathrm{SG})$ is def\/ined
by naturally extending these examples.
Namely, we consider $n-1$ quivers $Q_1,\dots, Q_{n-1}$,
whose vertices are naturally identif\/ied with
the vertices of the graph $X_n$ in Fig.~\ref{fig:X}.
(The leftmost vertex corresponds to the vertex~$1$ in~$X_n$.)
The arrows are put in $Q_i$ as clearly indicated by the examples
in Fig.~\ref{fig:quiverG}.
Note that the pattern of arrows slightly
depends on the parity of~$n$.
Then, except for the leftmost vertex of each quiver~$Q_i$,
 all the  vertices in the same position
 in  $n-1$ quivers $Q_1,\dots,Q_{n-1}$ are
identif\/ied.
Also we assign the property $+/-$
to  each vertex, except for the leftmost one in each~$Q_i$,
 as in Fig.~\ref{fig:quiverG}.

Let us choose  the index set $\mathbf{I}$
of the vertices of $Q_{n}(\mathrm{SG})$
so that $\mathbf{i}=(i,1)\in \mathbf{I}$ represents
the leftmost vertex in $Q_i$ for $i=1,\dots,n-1$,
 and $\mathbf{i}=(n,i')\in \mathbf{I}$ represents
the vertex $i'=2,\dots,n+1$ in any quiver $Q_i$
under the natural
identif\/ication with $X_n$.
Thus, $i=1,\dots,n$; and $i'=1$ if $i\neq n$
and $i'=2,\dots,n+1$ if $i=n$.

Let $\mathbf{I}_+$
(resp.\ $\mathbf{I}_-$)
denote the set of the vertices $\mathbf{i}\in \mathbf{I}$ with
property  $+$ (resp.\  $-$).
We def\/ine composite mutations (Appendix~\ref{sec:groc}(ii)--(v)),
\begin{gather*}
%\label{eq:mupm2}
\mu_+=\prod_{\mathbf{i}\in\mathbf{I}_+}
\mu_{\mathbf{i}},
\qquad
\mu_-=\prod_{\mathbf{i}\in\mathbf{I}_-}
\mu_{\mathbf{i}}.
\end{gather*}
Note that they do not depend on the order of the product.

For a permutation $\sigma$ of $\{1,\dots,n-1 \}$,
let $\tilde{\sigma}$ be the permutation of
$\mathbf{I}$ such that
$\tilde{\sigma}(i,1)=(\sigma(i),1)$ for $i\neq n$
and $\tilde{\sigma}(n,i')=(n,i')$.
Let $\tilde{\sigma}(Q_{n}(\mathrm{SG}))$
denote the
quiver induced from $Q_{n}(\mathrm{SG})$
 by~$\tilde{\sigma}$.
Namely, if there is an arrow
 $\mathbf{i}\rightarrow
\mathbf{j}$ in $Q_{n}(\mathrm{SG})$,
then, there is an arrow
$\tilde{\sigma}(\mathbf{i})
\rightarrow
\tilde{\sigma}(\mathbf{j})
$
in  $\tilde{\sigma}(Q_{n}(\mathrm{SG}))$.

\begin{Lemma}
\label{lem:GQmut}
Let $Q(0):=Q_{n}(\mathrm{SG})$.
We have the following periodic sequence of mutations of quivers:
\begin{gather}
% 1st
Q(0)
\mathop{\longleftrightarrow}^{{\mu_+} \mu_{(1,1)}}
 Q(1)
\mathop{\longleftrightarrow}^{\mu_-}
Q(2)
\displaystyle
\mathop{\longleftrightarrow}^{\mu_+ \mu_{(2,1)}}
Q(3)
\displaystyle
\mathop{\longleftrightarrow}^{\mu_-}
Q(4)
\nonumber\\
\hphantom{Q(0)
\mathop{\longleftrightarrow}^{{\mu_+} \mu_{(1,1)}}}{}
\mathop{\longleftrightarrow}^{\mu_+\mu_{(3,1)}}
\ \cdots\
\displaystyle
\mathop{\longleftrightarrow}^{\mu_+\mu_{(n-1,1)}}
Q(2n-3)
\displaystyle
\mathop{\longleftrightarrow}^{\mu_-}
Q(2n-2)=Q(0),\label{eq:GB2}
\end{gather}
where the quiver $Q(2p)$ $(p=1,\dots,n-2)$ is given by
\begin{gather}
\label{eq:Q2p}
Q(2p)=
\tilde{\sigma}^p(Q(0)),
\qquad
\sigma=\left(
\begin{matrix}
1&2&\dots &n-2 &n-1\\
2&3& \dots&n-1& 1\\
\end{matrix}
\right).
\end{gather}
\end{Lemma}

\begin{Example}
The mutation sequence
 \eqref{eq:GB2} for $n=6$ is explicitly given
in Fig.~\ref{fig:labelxG1}.
\begin{figure}[t]
\centerline{\includegraphics{Nakanishi-Fig3}}

\caption{The mutation sequence of the quiver $Q_{n}(\mathrm{SG})$
in~\eqref{eq:GB2} for $n=6$.
The encircled vertices correspond to
the mutation points
$(\mathbf{i},u): \mathbf{p}_+$ in the forward direction.}
\label{fig:labelxG1}
\end{figure}
\end{Example}



\subsection{Embedding maps}
Let
 $B=B_{n}(\mathrm{SG})$
be the skew-symmetric matrix corresponding
to the quiver $Q_{n}(\mathrm{SG})$ (Appendix~\ref{sec:groc}(iii)).
Let $\mathcal{A}(B,x,y)$
be the cluster algebra
 with coef\/f\/icients
in the universal
semif\/ield
$\mathbb{P}_{\mathrm{univ}}(y)$ (Appendix~\ref{sec:groc}(vi)).



In view of Lemma~\ref{lem:GQmut}
we set $x(0)=x$, $y(0)=y$ and def\/ine
clusters $x(u)=(x_{\mathbf{i}}(u))_{\mathbf{i}\in \mathbf{I}}$
 ($u\in \mathbb{Z}$)
 and coef\/f\/icient tuples $y(u)=(y_\mathbf{i}(u))_{\mathbf{i}\in \mathbf{I}}$
 ($u\in \mathbb{Z}$)
by the sequence of mutations
\begin{gather}
% 1st
\cdots
\ \mathop{\longleftrightarrow}^{\mu_-}\
(B(0),x(0),y(0))
\ \mathop{\longleftrightarrow}^{\mu_+
\mu_{(1,1)}}\
(B(1),x(1),y(1))\nonumber\\
\phantom{\cdots}{} \ \mathop{\longleftrightarrow}^{\mu_-}\
\cdots
\ \mathop{\longleftrightarrow}^{\mu_-}\
(B(2n-2),x(2n-2),y(2n-2))
\ \mathop{\longleftrightarrow}^{\mu_+
\mu_{(1,1)}}\
\cdots,\label{eq:Gmutseq}
\end{gather}
where $B(u)$ is the skew-symmetric matrix corresponding to
$Q(u)$.

For  $(\mathbf{i},u)\in
 \mathbf{I}\times \mathbb{Z}$,
we set the parity condition $\mathbf{p}_+$ by
\begin{gather}
\label{eq:GQparity}
\mathbf{p}_+: \
\begin{cases}
 \mathbf{i}\in
\mathbf{I}_+ \sqcup
\{(j+1,1)\},
& u\equiv 2j
\ (j=0,1,\dots,n-2),\\
 \mathbf{i}\in
\mathbf{I}_-,
& \mbox{$u$: odd},
\end{cases}
\end{gather}
where $\equiv$ is modulo $(2n-2)\mathbb{Z}$.
We def\/ine the condition $\mathbf{p}_-$
by $(\mathbf{i},u):\mathbf{p}_- \Longleftrightarrow
(\mathbf{i},u-1):\mathbf{p}_+$.
Plainly speaking,
each $(\mathbf{i},u):\mathbf{p}_+$
(resp.\ $(\mathbf{i},u):\mathbf{p}_-$)
is a mutation point of~\eqref{eq:Gmutseq} in the forward
(resp. backward) direction of~$u$.
See Fig.~\ref{fig:labelxG1}.

\begin{Lemma}
\label{lem:gmap}
Below $\equiv$ means the equivalence modulo $(2n-2)\mathbb{Z}$.

\begin{enumerate}\itemsep=0pt
\item[$(i)$]
The map
\begin{gather*}
g:  \ \mathcal{I}_{n+} \rightarrow
\{ (\mathbf{i},u)\in \mathbf{I}\times \mathbb{Z} \mid
(\mathbf{i},u): \mathbf{p}_+
\}, \\
\phantom{g: \ } \ (i,u-d_i) \mapsto
\begin{cases}
((j+1,1),u), & i= 1; \
u \equiv 2j
\ (j=0,1,\dots,n-2),\\
((n,i),u), & i=2,\dots,n+1
\end{cases}
\end{gather*}
is a bijection.

\item[$(ii)$]
The map
\begin{gather*}
g': \  \mathcal{I}'_{n+} \rightarrow
\{ (\mathbf{i},u)\in \mathbf{I}\times \mathbb{Z} \mid
 (\mathbf{i},u): \mathbf{p}_+
\}, \\
\phantom{g': \ }{} \  (i,u) \mapsto
\begin{cases}
((j+1,1),u), & i= 1;\
u\equiv 2j
\ (j=0,1,\dots,n-2), \\
((n,i),u), &  i=2,\dots,n+1
\end{cases}
\end{gather*}
is a bijection.
\end{enumerate}
\end{Lemma}

Based on Lemma \ref{lem:gmap},
we introduce alternative notations
$\tilde{x}_i(u-d_i):=x_{\mathbf{i}}(u)$
for $(i,u-d_i)\in \mathcal{I}_{n+}$
with $(\mathbf{i},u)=g((i,u-d_i))$
and
$y_i(u):=y_{\mathbf{i}}(u)$
for $(i,u)\in \mathcal{I}'_{n+}$
with $(\mathbf{i},u)=g'((i,u))$,
respectively,
which turn out to be useful.

 Let $\mathcal{A}(B,x)$ be the cluster algebra
with trivial coef\/f\/icients, where $(B,x)$ is
the initial seed
and the coef\/f\/icient semif\/ield
is the {\em trivial semifield\/}
 $\mathbf{1}=\{1\}$ (Appendix~\ref{sec:groc}(i)).
Let $\pi_{\mathbf{1}}:
\mathbb{P}_{\mathrm{univ}}(y)\rightarrow
\mathbf{1}$, $y_{\mathbf{i}}\mapsto 1$ be the projection.
Let $[x_{\mathbf{i}}(u)]_{\mathbf{1}}$
denote the image of $x_{\mathbf{i}}(u)$
 by the algebra homomorphism
$\mathcal{A}(B,x,y)\rightarrow \mathcal{A}(B,x)$
 induced from $\pi_{\mathbf{1}}$.
It is called the {\em trivial evaluation}.




The following lemma follows from the exchange relation
of
cluster variables \eqref{eq:clust}
and the property of the sequence \eqref{eq:GB2}
 observed in Fig.~\ref{fig:labelxG1}.

\begin{Lemma}
\label{lem:Gx2}
Let $G_+$ and $G_-$ be the ones in \eqref{eq:Yuni} and \eqref{eq:Tuni}.
The family $\{\tilde{x}_i(u)
\mid (i,u)\in \mathcal{I}_{n+}\}$
satisfies a system of relations
\begin{gather}
\tilde{x}_i\left(u-d_i\right)
\tilde{x}_i\left(u+d_i\right)
 =
\frac{y_i(u)}{1+y_i(u)}
\prod_{(j,v)\in \mathcal{I}_{n+}}
\tilde{x}_j(v)^{G_+(i,u;j,v)}
\nonumber\\
\phantom{\tilde{x}_i\left(u-d_i\right)
\tilde{x}_i\left(u+d_i\right)
 =}{}
+
\frac{1}{1+y_i(u)}
\prod_{(j,v)\in \mathcal{I}_{n+}}
\tilde{x}_j(v)^{G_-(i,u;j,v)},\label{eq:xy1}
\end{gather}
where $(i,u)\in \mathcal{I}'_{n+}$.
In particular,
the family $\{ [\tilde{x}_i(u)]_{\mathbf{1}}
\,|\, (i,u)\in \mathcal{I}_{n+}\}$
satisfies the T-system $\mathbb{T}_{n}(\mathrm{SG})$
in $\mathcal{A}(B,x)$
by replacing $T_i(u)$ with $[\tilde{x}_i(u)]_{\mathbf{1}}$.
\end{Lemma}

\begin{Example}
Consider the case $n=6$ in Fig.~\ref{fig:labelxG1}.
Let us consider the mutation at the vertex~$(1,1)$ in $Q(0)$,
to which  the variable $\tilde{x}_1(-5)$ is attached.
The next time~$(1,1)$ is mutated is in~$Q(10)$,
where~$\tilde{x}_1(5)$ is attached.
Meanwhile, the only vertex connected to~$(1,1)$ in~$Q(0)$ is~$(6,2)$,
and the  variable attached to $(6,2)$ in $Q(0)$ is equal to
the variable  $\tilde{x}_2(0)$ attached to~$(6,2)$ in~$Q(1)$.
Taking account of the directions of the arrows,
we have the  relation
\begin{align*}
\tilde{x}_1(-5)\tilde{x}_1(5)=
\frac{y_1(0)}{1+y_1(0)} \tilde{x}_2(0)
+ \frac{1}{1+y_1(0)} ,
\end{align*}
which agrees with
\eqref{eq:T1} and
\eqref{eq:xy1}.
Similarly, consider the mutation at the vertex $(6,2)$ in $Q(1)$,
to which  the variable $\tilde{x}_2(0)$ is attached.
The next time $(6,2)$ is mutated is in $Q(3)$,
where $\tilde{x}_2(2)$ is attached.
Meanwhile, the vertices connected to $(6,2)$ in $Q(1)$
are $(1,1)$, $(2,1)$, and $(6,3)$,
and the  variable attached to them are equal to
$\tilde{x}_1(5)$, $\tilde{x}_1(-3)$, and $\tilde{x}_3(1)$, respectively.
Taking account of the directions of the arrows,
we have the relation
\begin{align*}
\tilde{x}_2(0)\tilde{x}_2(2)=
\frac{y_2(1)}{1+y_2(1)} \tilde{x}_1(-3)\tilde{x}_1(5)
+ \frac{1}{1+y_2(1)} \tilde{x}_3(1).
\end{align*}
As the last example,
consider the mutation at the vertex $(6,6)$ in $Q(1)$,
to which  the va\-riab\-le~$\tilde{x}_6(0)$ is attached.
The next time $(6,6)$ is mutated is in $Q(3)$,
where $\tilde{x}_6(2)$ is attached.
Meanwhile, the vertices connected to $(6,6)$ in $Q(1)$
are $(4,1)$ and $(6,5)$,
and the  variable attached to them are equal to
$\tilde{x}_1(1)$ and  $\tilde{x}_5(1)$, respectively.
Taking account of the directions of the arrows,
we have the relation
\[
\tilde{x}_6(0)\tilde{x}_6(2)=
\frac{y_6(1)}{1+y_6(1)} \tilde{x}_1(1)
+ \frac{1}{1+y_6(1)} \tilde{x}_5(1).
\]
The other cases can be checked in similar manners.
\end{Example}

\begin{Definition}
\label{def:T-sub}
The {\em T-subalgebra
$\mathcal{A}_T(B,x)$
of ${\mathcal{A}}(B,x)$
associated with the sequence \eqref{eq:Gmutseq}}
is the subring of
${\mathcal{A}}(B,x)$
generated by
$[x_{\mathbf{i}}(u)]_{\mathbf{1}}$
($(\mathbf{i},u)\in \mathbf{I}\times \mathbb{Z}$),
or equivalently,
generated by
$[\tilde{x}_{i}(u)]_{\mathbf{1}}$
($(i,u)\in \mathcal{I}_{n+}$).
\end{Definition}

By Lemma \ref{lem:Gx2}, we have the following embedding.
\begin{Theorem}
\label{thm:GTiso}
The ring $\EuScript{T}^{\circ}_{n}(\mathrm{SG})_+$ is isomorphic to
$\mathcal{A}_T(B,x)$ by the correspondence
$T_i(u)\mapsto [\tilde{x}_i(u)]_{\mathbf{1}}$.
\end{Theorem}



The {\em coefficient group $\mathcal{G}(B,y)$
associated with $\mathcal{A}(B,x,y)$}
is the multiplicative subgroup of
the semif\/ield $\mathbb{P}_{\mathrm{univ}}(y)$ generated by all
the coef\/f\/icients $y_{\mathbf{i}}'$ of $\mathcal{A}(B,x,y)$
together with $1+y_{\mathbf{i}}'$.

The following lemma follows from the exchange relation of
coef\/f\/icients \eqref{eq:coef} and the property of the sequence
\eqref{eq:GB2}.


\begin{Lemma}
\label{lem:Gy2}
The family $\{ y_i(u)
\mid (i,u)\in \mathcal{I}'_{n+}\}$
satisfies the Y-system $\mathbb{Y}_{n}(\mathrm{SG})$
by replacing $Y_i(u)$ with $y_i(u)$.
\end{Lemma}

\begin{Example}
Consider the case $n=6$ in Fig.~\ref{fig:labelxG1}.
Let us consider the mutation at the vertex~$(1,1)$ in $Q(0)$,
to which  the variable $y_1(0)$ is attached.
The next time~$(1,1)$ is mutated is in~$Q(10)$,
where $y_1(10)$ is attached.
Meanwhile, between $u=0$ and $u=10$,
the vertices connected to $(1,1)$ in~$Q(u)$ and mutated
are
$(6,2)$ at $u=1,9$, $(6,3)$ at $u=2,8$, $(6,4)$ at $u=3,7$,
$(6,5)$ at $u=4,6$, and  $(6,6)$ and $(6,7)$ at~$u=5$,
Taking account of the directions of the arrows,
we have the relation
\begin{gather*}
y_1(0)y_1(10)=
(1+y_2(1))(1+y_2(9))(1+y_3(2))(1+y_3(8)) (1+y_4(3))\\
\phantom{y_1(0)y_1(10)=}{} \times (1+y_4(7))
(1+ y_5(4))(1+y_5(6)) (1+y_6(5))(1+y_7(5)),
\end{gather*}
which agrees with
\eqref{eq:Y1}.
Similarly, consider the mutation at the vertex $(6,2)$ in $Q(1)$,
to which  the variable $y_2(1)$ is attached.
The next time $(6,2)$ is mutated is in $Q(3)$,
where $y_2(3)$ is attached.
Meanwhile, between $u=1$ and $u=3$,
the vertices connected to $(6,2)$ in $Q(u)$ and mutated
are~$(2,1)$ and~$(6,3)$ at $u=2$.
Taking account of the directions of the arrows,
we have the relation
\[
y_2(1)y_2(3) =
\frac{(1+y_1(2))}{1+y_3(2)^{-1}}.
\]
As the last example, consider the mutation at the vertex $(6,6)$ in $Q(1)$,
to which  the variable~$y_6(1)$ is attached.
The next time $(6,6)$ is mutated is in $Q(3)$,
where $y_6(3)$ is attached.
Meanwhile, between $u=1$ and $u=3$,
the only vertex connected to $(6,6)$ in $Q(u)$ and mutated
is  $(6,5)$ at~$u=2$.
Taking account of the directions of the arrows,
we have the relation
\[
y_6(1)y_6(3) =
\frac{1}{1+y_5(2)^{-1}}.
\]
The other cases can be checked in similar manners.
\end{Example}

\begin{Definition}
The {\em Y-subgroup
$\mathcal{G}_Y(B,y)$
of $\mathcal{G}(B,y)$
associated with the sequence \eqref{eq:Gmutseq}}
is the subgroup of
$\mathcal{G}(B,y)$
generated by
$y_{\mathbf{i}}(u)$
($(\mathbf{i},u)\in \mathbf{I}\times \mathbb{Z}$)
and $1+y_{\mathbf{i}}(u)$
($(\mathbf{i},u):\mathbf{p}_+$ or $\mathbf{p}_-$),
or equivalently,
generated by
$y_i(u)$ and $1+y_i(u)$
($(i,u)\in \mathcal{I}'_{n+}$).
\end{Definition}

By Lemma \ref{lem:Gy2}, we have the following embedding.
\begin{Theorem}
\label{thm:GYiso}
The group $\EuScript{Y}^{\circ}_{n}(\mathrm{SG})_+$ is isomorphic to
$\mathcal{G}_Y(B,y)$ by the correspondence
$Y_i(u)\mapsto y_i(u)$
and $1+Y_i(u)\mapsto 1+y_i(u)$.
\end{Theorem}


\section{Proof of
Theorems \ref{thm:Yperiod}, \ref{thm:dilog},
and \ref{thm:Tperiod}}
\label{sec:proofSG}

In this section we prove
Theorems \ref{thm:Yperiod}, \ref{thm:dilog},
and \ref{thm:Tperiod}
using the method of \cite{Nakanishi09,Inoue10a}.


Let $y=y(0)$ be the initial coef\/f\/icient tuple
of the cluster algebra $\mathcal{A}(B,x,y)$
with $B=B_{n}(\mathrm{SG})$ in the previous section.
Let
$\mathbb{P}_{\mathrm{trop}}(y)$ be
the {\em tropical semifield} for $y$ (Appendix~\ref{sec:groc}(i)).
Let $\pi_{\mathbf{T}}:
\mathbb{P}_{\mathrm{univ}}(y)\rightarrow
\mathbb{P}_{\mathrm{trop}}(y)$, $y_{\mathbf{i}}\mapsto
y_{\mathbf{i}}$ be the projection.
Let $[y_{\mathbf{i}}(u)]_{\mathbf{T}}$
and $[\mathcal{G}_Y(B,y)]_{\mathbf{T}}$
denote the images of $y_{\mathbf{i}}(u)$
and $\mathcal{G}_Y(B,y)$
by the multiplicative group
 homomorphism induced from $\pi_{\mathbf{T}}$, respectively.
They are called the {\em tropical evaluations},
and the resulting relations in
the group $[\mathcal{G}_Y(B,y)]_{\mathbf{T}}$
are called the {\em tropical Y-system}.
They are f\/irst studied in~\cite{Fomin03b}
for cluster algebras of f\/inite types.

We say a (Laurent) monomial $m=\prod_{\mathbf{i}\in \mathbf{I}}
y_{\mathbf{i}}^{k_{\mathbf{i}}}$
is {\em positive} (resp. {\em negative})
if $m\neq 1$ and $k_{\mathbf{i}}\geq 0$
(resp. $k_{\mathbf{i}}\leq 0$)
for any $\mathbf{i}$.
It is known that every monomial $[y_{\mathbf{i}}(u)]_{\mathbf{T}}$
is either positive or negative by \cite[Proposition 5.6]{Fomin07}
and \cite[Theorem 1.7]{Derksen10}.

The next `tropical mutation rule' for $[y_{\mathbf{i}}(u)]_{\mathbf{T}}$
is general and  useful.
\begin{Lemma}
\label{lem:mutation}
Suppose that $y''$ is the coefficient tuple obtained from
the mutation of another coefficient tuple $y'$ at $\mathbf{k}$
with mutation matrix $B'$. Then, for any $\mathbf{i}\neq \mathbf{k}$,
we have the rule:
\begin{enumerate}\itemsep=0pt
\item[$(i)$] $[y''_{\mathbf{i}}]_{\mathbf{T}} =
 [y'_{\mathbf{i}}]_{\mathbf{T}}[y'_{\mathbf{k}}]_{\mathbf{T}}$
 if one of the following conditions holds:
\begin{enumerate}\itemsep=0pt
\item[$(a)$] $B'_{\mathbf{k}\mathbf{i}}> 0$,
 and $[y'_{\mathbf{k}}]_{\mathbf{T}}$ is
positive;
\item[$(b)$]  $B'_{\mathbf{k}\mathbf{i}}< 0$,
 and $[y'_{\mathbf{k}}]_{\mathbf{T}}$ is
negative.
\end{enumerate}
\item[$(ii)$]
$[y''_{\mathbf{i}}]_{\mathbf{T}} =
 [y'_{\mathbf{i}}]_{\mathbf{T}}$
 if one of the following conditions holds:
\begin{enumerate}\itemsep=0pt
\item[$(a)$] $B'_{\mathbf{k}\mathbf{i}}=0$;

\item[$(b)$]
 $B'_{\mathbf{k}\mathbf{i}}> 0$,
 and $[y'_{\mathbf{k}}]_{\mathbf{T}}$ is
 negative;
\item[$(c)$] $B'_{\mathbf{k}\mathbf{i}}< 0$,
 and $[y'_{\mathbf{k}}]_{\mathbf{T}}$ is
 positive.
 \end{enumerate}
 \end{enumerate}
\end{Lemma}
\begin{proof}
This is an immediate consequence of the exchange relation
\eqref{eq:coef} and \eqref{eq:trop}.
\end{proof}


The following properties
of the tropical Y-system are crucial.

\begin{Proposition}
\label{prop:lev2}
 For
$[\mathcal{G}_Y(B,y)]_{\mathbf{T}}$
with $B=B_{n}(\mathrm{SG})$, the following facts hold.

\begin{enumerate}\itemsep=0pt
\item[$(i)$]  For $0 \le u < 4n-2$,
the  monomial $[y_{\mathbf{i}}(u)]_{\mathbf{T}}$
$((\mathbf{i},u):\mathbf{p}_+)$
is negative if and only if $u$ takes the following values.
\[
\begin{cases}
2n-2 \leq u < 4n-2 &
\quad\mbox{\rm for $\mathbf{i}=(1,1),\dots,(n-1,1),(n,2)$},\\
u=2n-2,2n-1, 4n-4, 4n-3 &
\quad\mbox{\rm for $\mathbf{i}=(n,3),\dots,(n,n+1)$}.
\end{cases}
\]
$($Note that for each $\mathbf{i}$, $u$ takes only a part
of the list due to the condition $(\mathbf{i},u):\mathbf{p}_+.)$
\item[$(ii)$]
 We have $[y_{\mathbf{i}}(4n-2)]_{\mathbf{T}}
=y_{\tau^{-1}(\mathbf{i})}$,
where $\tau$ is a bijection $\mathbf{I} \rightarrow
\mathbf{I}$ defined by
\begin{gather*}
(i,1)   \mapsto (\sigma(i),1),
\qquad i=1,\dots,n-1,\\
(n,i')   \mapsto
\begin{cases}
(n,i'),& i'=2,\dots,n-1,\\
(n,n+1),& i'=n,\\
(n,n),& i'=n+1
\end{cases}
\end{gather*}
and $\sigma$ is the permutation in~\eqref{eq:Q2p}.

\item[$(iii)$] The number $N_-$ of the negative monomials
$[y_{\mathbf{i}}(u)]_{\mathbf{T}}$ for $(\mathbf{i},u):\mathbf{p}_+$
in the region $0\leq u < 4n-2$ is $4n-2$.
\end{enumerate}
\end{Proposition}



\begin{proof}
$(i)$ Let us factorize $[y_{\mathbf{i}}(u)]_{\mathbf{T}}
=[y_{\mathbf{i}}(u)]'_{\mathbf{T}}
[y_{\mathbf{i}}(u)]''_{\mathbf{T}}$,
where
 $[y_{\mathbf{i}}(u)]'_{\mathbf{T}}$ is a monomial
in $y_{(i,1)}$ ($i=1,\dots,n-1$)
while
 $[y_{\mathbf{i}}(u)]''_{\mathbf{T}}$ is a monomial
in $y_{(n,i')}$ ($i'=2,\dots,n+1$).
One can independently study
$[y_{\mathbf{i}}(u)]'_{\mathbf{T}}$ and
$[y_{\mathbf{i}}(u)]''_{\mathbf{T}}$.
The claim  (i) follows from the following results,
which  are proved inductively on $u$
by Lemma \ref{lem:mutation} and the results of
\cite{Fomin03b,Fomin07}.

$(a)$ $[y_{\mathbf{i}}(u)]'_{\mathbf{T}}$ part.
This part is easier.
All the monomials $[y_{\mathbf{i}}(u)]'_{\mathbf{T}}$
 which are not 1 for $(\mathbf{i},u):\mathbf{p}_+$
in the region $0 \leq u < 4n-2$ are as follows.
We have $[y_{(i,1)}(2i-2)]'_{\mathbf{T}}=y_{(i,1)}$
$(i=1,\dots,n-1)$,
and also
\begin{alignat}{3}
& [y_{(1,1)}(2n-2)]'_{\mathbf{T}}=y_{(1,1)}^{-1},\qquad &&
[y_{(n,2)}(2n-1)]'_{\mathbf{T}}=y_{(1,1)}^{-1},&\notag\\
& [y_{(2,1)}(2n)]'_{\mathbf{T}}=y_{(1,1)}^{-1}y_{(2,1)}^{-1},\qquad &&
[y_{(n,2)}(2n+1)]'_{\mathbf{T}}=y_{(2,1)}^{-1},&\notag\\
\label{eq:y'part}
& [y_{(3,1)}(2n+2)]'_{\mathbf{T}}=y_{(2,1)}^{-1}y_{(3,1)}^{-1},\qquad &&
[y_{(n,2)}(2n+3)]'_{\mathbf{T}}=y_{(3,1)}^{-1},&\\
&\qquad \vdots & &\qquad \vdots & \notag\\
& [y_{(n-1,1)}(4n-6)]'_{\mathbf{T}}=y_{(n-2,1)}^{-1}y_{(n-1,1)}^{-1},\qquad &&
[y_{(n,2)}(4n-5)]'_{\mathbf{T}}=y_{(n-1,1)}^{-1},&\notag\\
& [y_{(1,1)}(4n-4)]'_{\mathbf{T}}=y_{(n-1,1)}^{-1}.&&& \notag
\end{alignat}

$(b)$  $[y_{\mathbf{i}}(u)]''_{\mathbf{T}}$ part.
Below we list all the monomials $[y_{\mathbf{i}}(u)]''_{\mathbf{T}}$
which are not 1 for $(\mathbf{i},u):\mathbf{p}_+$
in the region $0 \leq u < 4n-2$.
We separate the region $0 \leq u < 4n-2$ into four parts
corresponding to the decomposition
$4n-2=(2n-2)+2+(2n-4)+2$,
where $2n-2$ and $2n-4$ are the Coxeter numbers
of $D_{n}$ and $D_{n-1}$.


Region I: $0\leq u < 2n-2$.
All the  monomials
$[y_{(n,i')}(u)]''_{\mathbf{T}}$ ($i'=2,\dots,n+1$)
 for $(\mathbf{(n,i')},u):\mathbf{p}_+$
are identif\/ied with
the {\em positive roots of $D_n$} as in
\cite[Proposition 10.7]{Fomin07};
therefore, they are positive.
Here, $D_{n}$ is identif\/ied with the subgraph of $X_n$
consisting of vertices $2, \dots, n+1$.


Region II: $u=2n-2$, $2n-1$.
We have, for even $n$,
\begin{gather*}
[y_{(n,i')}(2n-2)]''_{\mathbf{T}} =y_{(n,i')}^{-1},\qquad  i'=3,5,\dots,n-1 ,\nonumber\\
[y_{(n,i')}(2n-1)]''_{\mathbf{T}} =y_{(n,i')}^{-1},\qquad  i'=2,4,\dots,n,n+1 ,%\label{eq:yt1}
\end{gather*}
and, for odd $n$,
\begin{gather*}
[y_{(n,i')}(2n-2)]''_{\mathbf{T}} =y_{(n,i')}^{-1}, \qquad  i'=3,5,\dots,n-2 ,\nonumber\\
[y_{(n,n)}(2n-2)]''_{\mathbf{T}} =y_{(n,n+1)}^{-1},
\qquad [y_{(n,n+1)}(2n-2)]''_{\mathbf{T}}=y_{(n,n)}^{-1},\nonumber\\
[y_{(n,i')}(2n-1)]''_{\mathbf{T}} =y_{(n,i')}^{-1},\qquad i'=2,4,\dots,n-1.%\label{eq:yt4}
\end{gather*}

Region III: $2n \leq u < 4n-4$.
All the  monomials
$[y_{(n,i')} (u)]''_{\mathbf{T}}$ ($i'=3,\dots,n+1$)
 for $((n,i'),u):\mathbf{p}_+$
are  identif\/ied with
the {\em positive roots of $D_{n-1}$},
therefore, they are positive.
Here, $D_{n-1}$ is identif\/ied with the subgraph of $X_n$
consisting of vertices $3, \dots, n+1$.

Region IV: $u=4n-4$, $4n-3$. We have, for even $n$,
\begin{gather}
[y_{(n,i')}(4n-4)]''_{\mathbf{T}} =y_{(n,i')}^{-1}, \qquad i'=3,5,\dots,n-1,\nonumber\\
[y_{(n,i')}(4n-3)]''_{\mathbf{T}} =y_{(n,i')}^{-1}, \qquad i'=2,4,\dots,n-2,\nonumber\\
[y_{(n,n)}(4n-3)]''_{\mathbf{T}} =y_{(n,n+1)}^{-1},
\qquad [y_{(n,n+1)}(4n-3)]''_{\mathbf{T}}=y_{(n,n)}^{-1},\label{eq:yt3}
\end{gather}
and, for odd $n$,
\begin{gather}
[y_{(n,i')}(4n-4)]''_{\mathbf{T}} =y_{(n,i')}^{-1}, \qquad  i'=3,5,\dots,n-2 ,\nonumber\\
[y_{(n,n)}(4n-4)]''_{\mathbf{T}} =y_{(n,n+1)}^{-1},
\qquad [y_{(n,n+1)}(4n-4)]''_{\mathbf{T}}=y_{(n,n)}^{-1},\nonumber\\
[y_{(n,i')}(4n-3)]''_{\mathbf{T}} =y_{(n,i')}^{-1},\qquad i'=2,4,\dots,n-1.\label{eq:yt6}
\end{gather}

Besides, we have the sequences of  monomials
which appear over Regions~III and~IV;
for even~$n$,
\begin{gather}
 [y_{(2,1)}(2n)]''_{\mathbf{T}}=y_{(n,2)}^{-1}y_{(n,3)}^{-1},
\qquad [y_{(3,1)}(2n+2)]''_{\mathbf{T}}=y_{(n,4)}^{-1}y_{(n,5)}^{-1},
\qquad \dots,\nonumber\\
 [y_{(n/2,1)}(3n-4)]''_{\mathbf{T}}=y_{(n,n-2)}^{-1}y_{(n,n-1)}^{-1},\qquad
 [y_{(n/2+1,1)}(3n-2)]''_{\mathbf{T}}=y_{(n,n-1)}^{-1}y_{(n,n)}^{-1}y_{(n,n+1)}^{-1},\nonumber\\
 [y_{(n/2+2,1)}(3n)]''_{\mathbf{T}}=y_{(n,n-3)}^{-1}y_{(n,n-2)}^{-1},\qquad
\dots,\nonumber\\
  [y_{(n-1,1)}(4n-6)]''_{\mathbf{T}}=y_{(n,3)}^{-1}y_{(n,4)}^{-1},
\qquad [y_{(1,1)}(4n-4)]''_{\mathbf{T}}=y_{(n,2)}^{-1},\label{eq:yt2}
\end{gather}
and, for odd $n$, the middle three terms are replaced with
\begin{gather}
 [y_{((n-1)/2,1)}(3n-5)]''_{\mathbf{T}}=y_{(n,n-3)}^{-1}y_{(n,n-2)}^{-1},\nonumber\\
 [y_{((n+1)/2,1)}(3n-3)]''_{\mathbf{T}}=
y_{(n,n-1)}^{-1}y_{(n,n)}^{-1}y_{(n,n+1)}^{-1},\nonumber\\
 [y_{(n+3)/2,1)}(3n-1)]''_{\mathbf{T}}=y_{(n,n-2)}^{-1}y_{(n,n-1)}^{-1}.\label{eq:yt5}
\end{gather}

$(ii)$ They  follow from
\eqref{eq:yt3}--\eqref{eq:yt5}.


$(iii)$ By $(i)$, for each  $\mathbf{i}$
the numbers of the negative monomials
$[y_{\mathbf{i}}(u)]_{\mathrm{T}}$
in the region is~2 for $\mathbf{i}= (1,1)$,
 1 for $\mathbf{i}= (i,1)$ ($i=2,\dots,n-1$),
$n$ for $\mathbf{i}= (n,2)$.
and $2$ for $\mathbf{i}=(n,i')$ ($i'=3,\dots,n+1$).
Summing up, we have $N_-=4n-2$.
\end{proof}

Now we prove
Theorems \ref{thm:Yperiod} and \ref{thm:Tperiod}.
It follows from  a very general theorem
\cite[Theorem 5.1]{Inoue10a}
(based on the work by Plamondon \cite{Plamondon10a,Plamondon10b})
that the cluster variables
$x_{\mathbf{i}}(u)$ and coef\/f\/icients $y_{\mathbf{i}}(u)$
have the same periodicity with
$[y_{\mathbf{i}}(u)]_{\mathbf{T}}$ as in Proposition \ref{prop:lev2}$(ii)$,
namely,
$x_{\mathbf{i}}(4n-2)=x_{\tau^{-1}(\mathbf{i})}$
and
$y_{\mathbf{i}}(4n-2)=y_{\tau^{-1}(\mathbf{i})}$.
It follows that, under the labelling introduced in Lemma
\ref{lem:gmap}, we have
\begin{gather}
\tilde{x}_i(u+4n-2) =\tilde{x}_{\omega(i)}(u),
\qquad (i,u):\mathbf{P}_+,\nonumber\\
y_i(u+4n-2) =y_{\omega(i)}(u),
\qquad (i,u):\mathbf{P}'_+,\label{eq:xperiod}
\end{gather}
where $\omega$ is the one in Theorem~\ref{thm:Yperiod}$(i)$.
Then, thanks to the isomorphisms
in Theorems~\ref{thm:GTiso} and~\ref{thm:GYiso},
and also by the isomorphisms
$\EuScript{T}^{\circ}_{n}(\mathrm{SG})_+
\simeq
\EuScript{T}^{\circ}_{n}(\mathrm{SG})_-
$
and
$\EuScript{Y}^{\circ}_{n}(\mathrm{SG})_+
\simeq
\EuScript{Y}^{\circ}_{n}(\mathrm{SG})_-
$,
we obtain
Theorems~\ref{thm:Yperiod} and~\ref{thm:Tperiod}.
